% GENERATED FILE. Edit canonical lessons, then run npm run generate:books.
% canonical-source-hash: fnv1a64:b5b908c2c491eb24
% canonical-lessons: FR-C50-cou, FR-C50-epaule, FR-C50-genou, FR-C50-visage, FR-C50-sang

\chapter{Neck, Shoulder, Knee, Face, Blood}
\label{ch:fr-neck-shoulder-knee-face-blood}
\hlchaptermodality{50}

\begin{chapteropening}
\textbf{By the end of this chapter:} \emph{I can name the neck, a shoulder, a knee, the face, blood.}

Complete the last lesson of chapter 50: I can name the neck, a shoulder, a knee, the face, blood.
\end{chapteropening}

\section[le cou]{le cou — the neck}

\label{lesson:FR-C50-cou}



\begin{quote}
Before the new one: say the French for a finger, then the French for a tooth.
\end{quote}

\subsection*{You'll want to know: le cou}
\textbf{le cou} — \textquotedblleft{}the neck\textquotedblright{}.

\begin{cousinweb}
\textbf{cou} $\leftarrow$ Latin \emph{collum}, \textquotedblleft{}neck\textquotedblright{} — inherited. (EtymDB 2.1) Old French said \emph{col}; the \emph{l} later melted into a vowel.
\end{cousinweb}

\begin{grammarlens}[title={What you've built}]
One more word for this chapter.
\end{grammarlens}

\subsection*{Your turn}
\begin{itemize}
\raggedright
  \item \emph{Say it:} \emph{le cou}
  \item \emph{Say it:} \emph{le cou}, once more
  \item \emph{Say it:} say \emph{le cou}
  \item \emph{Recall:} say the French for a finger, then the French for a tooth, then say \emph{le cou} again
\end{itemize}

\begin{tcolorbox}[breakable,colback=teal!4,colframe=teal!35!black,title={Before you move on}]
What does le cou mean? (\textquotedblleft{}The neck\textquotedblright{}.) Say it once more.
\end{tcolorbox}
\section[l'épaule]{l'épaule — a shoulder}

\label{lesson:FR-C50-epaule}



\begin{quote}
Before the new one: say the French for a tooth, then the French for the neck.
\end{quote}

\subsection*{You'll want to know: l'épaule}
\textbf{l'épaule} — \textquotedblleft{}a shoulder\textquotedblright{}.

\begin{cousinweb}
\textbf{épaule} $\leftarrow$ Old French \emph{espalle} $\leftarrow$ Late Latin \emph{spatula}, \textquotedblleft{}a flat piece, a shoulder blade\textquotedblright{} — inherited. (EtymDB 2.1) The \emph{é-} stands where an \emph{s} once was: \emph{espaule}, then \emph{épaule}.
\end{cousinweb}

\begin{grammarlens}[title={What you've built}]
One more word for this chapter.
\end{grammarlens}

\subsection*{Your turn}
\begin{itemize}
\raggedright
  \item \emph{Say it:} \emph{l'épaule}
  \item \emph{Say it:} \emph{l'épaule}, once more
  \item \emph{Say it:} say \emph{l'épaule}
  \item \emph{Recall:} say the French for a tooth, then the French for the neck, then say \emph{l'épaule} again
\end{itemize}

\begin{tcolorbox}[breakable,colback=teal!4,colframe=teal!35!black,title={Before you move on}]
What does l'épaule mean? (\textquotedblleft{}A shoulder\textquotedblright{}.) Say it once more.
\end{tcolorbox}
\section[le genou]{le genou — a knee}

\label{lesson:FR-C50-genou}



\begin{quote}
Before the new one: say the French for the neck, then the French for a shoulder.
\end{quote}

\subsection*{You'll want to know: le genou}
\textbf{le genou} — \textquotedblleft{}a knee\textquotedblright{}.

\begin{cousinweb}
\textbf{genou} $\leftarrow$ Vulgar Latin \emph{*genuculum}, \textquotedblleft{}a little knee\textquotedblright{}, from Latin \emph{genu}, \textquotedblleft{}knee\textquotedblright{} — inherited. (EtymDB 2.1 records both Latin words.)
\end{cousinweb}

\begin{grammarlens}[title={What you've built}]
One more word for this chapter.
\end{grammarlens}

\subsection*{Your turn}
\begin{itemize}
\raggedright
  \item \emph{Say it:} \emph{le genou}
  \item \emph{Say it:} \emph{le genou}, once more
  \item \emph{Say it:} say \emph{le genou}
  \item \emph{Recall:} say the French for the neck, then the French for a shoulder, then say \emph{le genou} again
\end{itemize}

\begin{tcolorbox}[breakable,colback=teal!4,colframe=teal!35!black,title={Before you move on}]
What does le genou mean? (\textquotedblleft{}A knee\textquotedblright{}.) Say it once more.
\end{tcolorbox}
\section[le visage]{le visage — the face}

\label{lesson:FR-C50-visage}



\begin{quote}
Before the new one: say the French for a shoulder, then the French for a knee.
\end{quote}

\subsection*{You'll want to know: le visage}
\textbf{le visage} — \textquotedblleft{}the face\textquotedblright{}.

\begin{cousinweb}
\textbf{visage} — built in Old French from \emph{vis}, \textquotedblleft{}face\textquotedblright{}, with the ending \emph{-age}. \emph{vis} $\leftarrow$ Latin \emph{visus}, \textquotedblleft{}a look, an appearance\textquotedblright{}, from \emph{videre}, \textquotedblleft{}to see\textquotedblright{}. (EtymDB 2.1)
\end{cousinweb}

\begin{grammarlens}[title={What you've built}]
One more word for this chapter.
\end{grammarlens}

\subsection*{Your turn}
\begin{itemize}
\raggedright
  \item \emph{Say it:} \emph{le visage}
  \item \emph{Say it:} \emph{le visage}, once more
  \item \emph{Say it:} say \emph{le visage}
  \item \emph{Recall:} say the French for a shoulder, then the French for a knee, then say \emph{le visage} again
\end{itemize}

\begin{tcolorbox}[breakable,colback=teal!4,colframe=teal!35!black,title={Before you move on}]
What does le visage mean? (\textquotedblleft{}The face\textquotedblright{}.) Say it once more.
\end{tcolorbox}
\section[le sang]{le sang — blood}

\label{lesson:FR-C50-sang}



\begin{quote}
Before the new one: say the French for a knee, then the French for the face.
\end{quote}

\subsection*{You'll want to know: le sang}
\textbf{le sang} — \textquotedblleft{}blood\textquotedblright{}.

\begin{cousinweb}
\textbf{sang} $\leftarrow$ Old French \emph{sanc} $\leftarrow$ Latin \emph{sanguis}, \textquotedblleft{}blood\textquotedblright{}, through a Vulgar Latin form \emph{*sanguem} — inherited. (EtymDB 2.1)
\end{cousinweb}

\begin{grammarlens}[title={What you've built}]
One more word for this chapter.
\end{grammarlens}

\subsection*{Your turn}
\begin{itemize}
\raggedright
  \item \emph{Say it:} \emph{le sang}
  \item \emph{Say it:} \emph{le sang}, once more
  \item \emph{Say it:} say \emph{le sang}
  \item \emph{Recall:} say the French for a knee, then the French for the face, then say \emph{le sang} again
\end{itemize}

\begin{tcolorbox}[breakable,colback=teal!4,colframe=teal!35!black,title={Before you move on}]
What does le sang mean? (\textquotedblleft{}Blood\textquotedblright{}.) Say it once more.
\end{tcolorbox}
